We now give a direct treatment of the classical and modified Yang–Baxter equations in the present framework \cite{SemenovTyanShanskii1983,Sklyanin1982}.

The $r$-matrix formalism can be formulated in a natural way in the setting of Lie algebras. In particular, the classical Yang--Baxter equation is most conveniently expressed in terms of tensor products of a Lie algebra and its invariant structures.

We therefore introduce the general algebraic framework that will be used throughout this section.

\begin{definition}
    Let $\mathcal{G}$ be a finite-dimensional Lie algebra over $\mathbb{C}$, equipped with an invariant non-degenerate bilinear form $(\cdot,\cdot)$. We denote by $\mathcal{G}^{\otimes n}$ the $n$-fold tensor product of $\mathcal{G}$, and use the standard notation $X_1, X_2, X_3$ for the embedding of $X \in \mathcal{G}$ into the corresponding tensor factors.
\end{definition}

\begin{definition}\label{def:double_bracket}
    Let $r \in \mathcal{G}^{\otimes 2}$. We define
    \begin{equation}
        [[r,r]] \equiv [r_{12},r_{13}] + [r_{12},r_{23}] + [r_{13},r_{23}]
        \in \mathcal{G}^{\otimes 3}.
    \end{equation}
\end{definition}
\begin{definition}
    Let $C \in \mathcal{G}^{\otimes 2}$ be the Casimir tensor element.  
    We define
    \begin{equation}
        \Omega \equiv [[C,C]] \in \mathcal{G}^{\otimes 3}.
    \end{equation}
\end{definition}
\begin{remark}
    The tensor $\Omega$ is called the canonical invariant element in $\mathcal{G}^{\otimes 3}$.
\end{remark}
\begin{definition}
    Let $r \in \mathcal{G}^{\otimes 2}$. The \emph{classical Yang--Baxter equation} (CYBE) is
    \begin{equation} \label{eq:CYBE}
        [[r,r]] = 0,
    \end{equation}
    while the \emph{modified classical Yang--Baxter equation} (mCYBE) is
    \begin{equation} \label{eq:mCYBE}
        [[r,r]] = \mu \, \Omega,
    \end{equation}
    where $\mu \in \mathbb{C}$ is a constant.
\end{definition}

\begin{proposition}
    Let $L$ be a Lax matrix whose Poisson brackets satisfy \eqref{eq:fund}, and assume that $r \in \mathcal{G}^{\otimes 2}$ is an antisymmetric constant tensor. If $r$ satisfies the modified classical Yang--Baxter equation, then the Jacobi identity for the Poisson bracket is satisfied:
    \begin{equation}
        \{\{L_1,L_2\},L_3\}
        + \{\{L_2,L_3\},L_1\}
        + \{\{L_3,L_1\},L_2\}
        = 0.
    \end{equation}
\end{proposition}
\begin{proof}
    We directly compute the terms in the Jacobi identity:
    \begin{equation} \label{eq:L_1 L_2 L_3}
        \{ \{L_1 , L_2 \} , L_3 \} \overset{(\ref{eq:fund})}{=} \{ [r_{12}, L_1] - [r_{21}, L_2] , L_3 \} = \{ [r_{12}, L_1] , L_3 \} - \{ [r_{21}, L_2] , L_3 \}.
    \end{equation}
    Since $r$ is assumed to be constant, it Poisson-commutes with $L$, and since it is antisymmetric we have
    \begin{equation}
        r_{21} = -r_{12}, \qquad r_{32} = -r_{23}, \qquad r_{31} = -r_{13}.
    \end{equation}
    We focus on 
    \begin{equation}
        \{ [r_{12}, L_1] , L_3 \} = \{ r_{12} L_1 - L_1 r_{12}, L_3 \} = r_{12} \{ L_1, L_3 \} - \{ L_1, L_3 \} r_{12} = [r_{12}, \{ L_1, L_3 \}],
    \end{equation}
    but $\{ L_1, L_3 \} = [r_{13}, L_1 ] - [r_{31}, L_3 ]$ then 
    \begin{equation}
        \{ [r_{12}, L_1] , L_3 \} = [r_{12}, [r_{13}, L_1]] - [r_{12}, [r_{31}, L_3]],
    \end{equation}
    in the same way we can compute the other term
    \begin{equation}
        \{ [r_{21}, L_2] , L_3 \} = [r_{21}, [r_{23}, L_2]] - [r_{21}, [r_{32}, L_3]].
    \end{equation}
    Hence \eqref{eq:L_1 L_2 L_3} reduces to the sum of 
    \begin{align} 
        \{ \{L_1, L_2 \}, L_3 \} &= [r_{12}, [r_{13}, L_1]] - [r_{12}, [r_{31}, L_3]] - [r_{21}, [r_{23}, L_2]] + [r_{21}, [r_{32}, L_3]] \label{eq:L_123},\\
        \{ \{L_3, L_1 \}, L_2 \} &= [r_{31}, [r_{32}, L_3]] - [r_{31}, [r_{23}, L_2]] - [r_{13}, [r_{12}, L_1]] + [r_{13}, [r_{21}, L_2]] \label{eq:L_321},\\
        \{ \{L_2, L_3 \}, L_1 \} &= [r_{23}, [r_{21}, L_2]] - [r_{23}, [r_{12}, L_1]] - [r_{32}, [r_{31}, L_3]] + [r_{32}, [r_{13}, L_1]] \label{eq:L_231}.
    \end{align}
    Summing \eqref{eq:L_123}, \eqref{eq:L_321}, and \eqref{eq:L_231}, we collect separately the terms proportional to $L_1$, $L_2$, and $L_3$.
    The terms containing $L_1$ are
    \begin{equation}
        [r_{12},[r_{13},L_1]] - [r_{13},[r_{12},L_1]] - [r_{23},[r_{12},L_1]] + [r_{32},[r_{13},L_1]].
    \end{equation}
    Using the Jacobi identity in the form $[a,[b,X]] - [b,[a,X]] = [[a,b],X]$  the first two terms give
    \begin{equation}
        [[r_{12},r_{13}],L_1].
    \end{equation}
    While for the last two terms we use Jacobi again to obtain
    \begin{equation}
        - [r_{23},[r_{12},L_1]] =
        - [[r_{23},r_{12}],L_1] - [r_{12},[r_{23},L_1]],
    \end{equation}
    and
    \begin{equation}
        [r_{32},[r_{13},L_1]] =
        [[r_{32},r_{13}],L_1] + [r_{13},[r_{32},L_1]].
    \end{equation}
    Hence the full contribution proportional to $L_1$ becomes
    \begin{equation} \label{eq:full_L1}
     [[r_{12},r_{13}],L_1] + [[r_{12},r_{23}],L_1] + [[r_{13},r_{23}],L_1] - [r_{12},[r_{23},L_1]] + [r_{13},[r_{32},L_1]].
    \end{equation}
    Since $r_{23}$ and $r_{32}$ act on the second and third tensor factors, while $L_1$ acts on the first one, we have
    \begin{equation}
        [r_{23},L_1] = [r_{32},L_1] = 0 \implies - [r_{12},[r_{23},L_1]] = 0 = [r_{13},[r_{32},L_1]],
    \end{equation}
    hence the residual terms vanish.
    Thus the full contribution proportional to $L_1$ becomes
    \begin{equation}
        [[r_{12},r_{13}] + [r_{12},r_{23}] + [r_{13},r_{23}], L_1].
    \end{equation}
    Using the modified classical Yang--Baxter equation, this is equal to
    \begin{equation}
        \mu [\Omega, L_1].
    \end{equation}
    Summing the three contributions, we obtain
    \begin{equation}
        \mu \Big( [\Omega,L_1] + [\Omega,L_2] + [\Omega,L_3] \Big)
        = \mu [\Omega, L_1 + L_2 + L_3].
    \end{equation}
    Now, by definition,
    \begin{equation}
        \Omega = [C_{12},C_{13}] + [C_{12},C_{23}] + [C_{13},C_{23}].
    \end{equation}
    
    Each of the tensors $C_{12}$, $C_{13}$ and $C_{23}$ commutes with $L_1+L_2+L_3$. 
    Indeed, by invariance of the Casimir tensor element, $C_{12}$ commutes with $L_1+L_2$, and it trivially commutes with $L_3$; similarly, $C_{13}$ commutes with $L_1+L_3$, and $C_{23}$ commutes with $L_2+L_3$. Therefore
    \begin{equation}
        [C_{12},L_1+L_2+L_3]=[C_{13},L_1+L_2+L_3]=[C_{23},L_1+L_2+L_3]=0.
    \end{equation}
    It follows that each commutator appearing in $\Omega$ also commutes with $L_1+L_2+L_3$, hence
    \begin{equation}
        [\Omega, L_1+L_2+L_3] = 0.
    \end{equation}
\end{proof}

\begin{lemma}
    Let $r_{12} \in \mathcal{G}^{\otimes2}$ be an antisymmetric solution of the modified classical Yang--Baxter equation, i.e.
    \begin{equation}
        r_{21} = -r_{12}, \qquad [[r,r]] = \mu \, \Omega .
    \end{equation}
    Let $C_{12}$ be the Casimir tensor element. For a choice of square root of $-\mu$, define
    \begin{equation}
        r_{12}^{\pm} = r_{12} \pm \sqrt{-\mu}\, C_{12}.
    \end{equation}
    Then $r_{12}^{\pm}$ are solutions of the classical Yang--Baxter equation:
    \begin{equation}\label{eq:cyb_pm}
        [r_{12}^{\pm},r_{13}^{\pm}] + [r_{12}^{\pm},r_{23}^{\pm}] + [r_{13}^{\pm},r_{23}^{\pm}] = 0.
    \end{equation}
\end{lemma}
\begin{proof}
    Substituting $r_{12}^{\pm}=r_{12}\pm \sqrt{-\mu} \, C_{12}$ into \eqref{eq:cyb_pm} and expanding, we obtain
    \begin{align}
        [r_{12},r_{13}] + [r_{12},&r_{23}] + [r_{13},r_{23}] \notag \\
        \pm \sqrt{-\mu} \bigg( 
        [C_{12},r_{13}] + [C_{12},r_{23}]
        + [r_{12},C_{13}] +& [r_{12},C_{23}]
        + [C_{13},r_{23}] + [r_{13},C_{23}]
        \bigg) \notag \\
        -\mu \bigg(
        [C_{12},C_{13}] + [C_{12},&C_{23}] + [C_{13},C_{23}]
        \bigg)
    \end{align}
    By \eqref{eq:casimir_YB}, the second term vanishes. Moreover, since $r_{12}$ satisfies the modified classical Yang--Baxter equation \eqref{eq:mCYBE}, the statement follows.
\end{proof}
